% ============================================================
% Section 9 — Relation to existing theories
% ============================================================

\section{Relation to existing theories}\label{sec:literature}

TGP occupies a distinctive niche among extensions of GR.
This section positions the theory relative to established
frameworks and identifies points of contact and departure.

\subsection{Scalar-tensor theories}

The closest relatives of TGP are Brans--Dicke theory and its
generalizations~\cite{Fujii2003,Damour1992}.  In these theories
a scalar field $\varphi$ is coupled to the Ricci scalar through
$\omega(\varphi) R$, and the PPN parameters depend on $\omega$.
TGP differs in three structural ways:
\begin{enumerate}[nosep]
  \item The field $\Phi$ is not a conformal factor of an
        \emph{a priori} metric; rather, the metric is an
        \emph{emergent} object constructed from $\Phi$.
  \item The kinetic coupling $K(\varphi)\propto\varphi^{4}$
        (with $\alpha=2$) is not a free function but is
        determined by the uniqueness theorem for $\mathcal{D}$
        and the vacuum condition $\beta=\gamma$.
  \item The fundamental constants $c$, $\hbar$, $G$ are all
        field-dependent, with $\ell_{P}=\mathrm{const}$ as the
        sole invariant.  This is more radical than VSL
        (variable speed of light) or VCC (varying coupling
        constant) models, yet avoids EEP violations because
        all dimensionless observables remain
        $\Phi$-independent~\cite{Uzan2011}.
\end{enumerate}

Despite these differences, TGP reproduces the GR limit with
$\gamma_{\mathrm{PPN}}=\beta_{\mathrm{PPN}}=1$
exactly~\cite{Will2014}, placing it in the same observational
equivalence class as GR in the weak-field regime.

\subsection{Gravitational wave constraints}

The observation of GW170817~\cite{GW170817} with its
electromagnetic counterpart constrains $|c_{T}/c-1|
\lesssim 10^{-15}$~\cite{Baker2017}, ruling out large classes
of scalar-tensor and vector-tensor models.  In TGP,
$c_{\mathrm{GW}}=c_{0}$ is an algebraic identity (the tensor
perturbation propagates on the substrate, not on $g_{\mu\nu}^{\mathrm{eff}}$),
so this constraint is satisfied trivially.

\subsection{Cosmological context}

The $\Lambda_{\mathrm{eff}}$ mechanism in TGP is qualitatively
similar to quintessence: a slowly evolving scalar potential
provides an effective equation of state $w(z)\neq -1$.
However, the physical origin is different: $\Lambda_{\mathrm{eff}}$
arises from \emph{fluctuations of generated space amplified
by matter condensation}, not from a fundamental potential.
The likelihood analysis against Planck~2018~\cite{Planck2018VI}
and DESI~DR1~\cite{DESI2024} data gives
$\Delta\chi^{2}=+1.8$, $\Delta\mathrm{AIC}=+3.8$ relative
to $\Lambda$CDM --- statistically indistinguishable.

\subsection{Fuzzy dark matter analogy}

The screened-Yukawa sector of TGP, with a Compton wavelength
$\lambda_{\mathrm{sp}}\sim\mathrm{Gpc}$, shares mathematical
structure with fuzzy dark matter
(FDM)~\cite{Hu2000,Schive2014}.  However, the effective boson
mass in TGP ($m_{\mathrm{sp}}\sim H_{0}$) is much larger
than the typical FDM candidate ($m\sim 10^{-22}\,\mathrm{eV}$).
Current Lyman-$\alpha$ constraints~\cite{Irsic2017,Rogers2021}
do not exclude the TGP regime; future DESI Lyman-$\alpha$ data
will provide a direct test (kill-shot K14).

\subsection{Black hole thermodynamics}

TGP reproduces the Bekenstein--Hawking
entropy~\cite{Bekenstein1973,Hawking1975} exactly:
$S_{\mathrm{BH}}=k_{B}A/(4\ell_{P}^{2})$, with the
$\Phi$-dependence cancelling due to
$\ell_{P}=\mathrm{const}$.  This is consistent with the
Wald entropy~\cite{Wald1993} in the weak-field limit where
the effective action reduces to Einstein--Hilbert.
In the strong-field regime (BH shadow, singularity
resolution), TGP makes predictions that differ from GR
and are testable by the Event Horizon
Telescope~\cite{EHT-M87-2019,EHT-SgrA-2022}.

\subsection{Particle sector: Koide formula}

The Koide relation $Q=(\sqrt{m_e}+\sqrt{m_\mu}+\sqrt{m_\tau})^2
/(m_e+m_\mu+m_\tau)=2/3$~\cite{Koide2007} has been known as
an empirical curiosity since 1982.  In TGP, it arises as a
\emph{consequence} of the soliton fixed-point structure,
not as an input.  The single parameter $g_{0}^{e}$
simultaneously determines lepton mass ratios (to $0.001\%$),
the Koide angle, and $\alpha_{s}(M_{Z})$~\cite{PDG2024}.
This is, to our knowledge, the first framework that derives
$Q=2/3$ from a dynamical mechanism rather than an
\emph{ad hoc} symmetry.

\subsection{Reproducibility}

All numerical results presented in this summary and in the
companion papers are reproducible from the open-source
repository~\cite{TGP-repo}, which contains 59 automated tests,
a 10-stage master reproduction pipeline, and GitHub Actions CI
running on Python~3.10/3.12.
